% =============================================================================
% embedded-glossar.tex – Glossar Embedded-Entwicklung
% =============================================================================

\section{Glossar: Embedded-Entwicklung}

Begriffe und Abkürzungen für dieses Projekt (ESP32-S3, CD4021, 74HC595, PlatformIO).

% -----------------------------------------------------------------------------
\subsection*{A}

\begin{description}[style=nextline, leftmargin=2.5cm, labelwidth=2.3cm]

\item[Active-Low]
Signal-Konvention: Aktiv = LOW (0), Inaktiv = HIGH (1).\\
\textit{Im Projekt:} CD4021-Buttons sind active-low (gedrückt = 0).

\item[ACK]
Acknowledge – Bestätigung einer empfangenen Nachricht.\\
\textit{Im Projekt:} Bridge sendet ACK nach erfolgreichem Kommando.

\end{description}

% -----------------------------------------------------------------------------
\subsection*{B}

\begin{description}[style=nextline, leftmargin=2.5cm, labelwidth=2.3cm]

\item[Baudrate]
Übertragungsgeschwindigkeit in Bit/s.\\
\textit{Im Projekt:} USB-CDC, Baudrate wird ignoriert (921600 nominal).

\item[Bit-Mapping]
Zuordnung von IDs zu Bit-Positionen.\\
\textit{Im Projekt:} MSB-first für Buttons und LEDs.

\end{description}

% -----------------------------------------------------------------------------
\subsection*{C}

\begin{description}[style=nextline, leftmargin=2.5cm, labelwidth=2.3cm]

\item[CD4021]
8-Bit Parallel-In/Serial-Out Schieberegister.\\
\textit{Im Projekt:} Liest Button-Zustände (Active-Low).

\item[COBS]
Consistent Overhead Byte Stuffing – Framing-Verfahren.\\
\textit{Im Projekt:} Für Serial-Protokoll zwischen ESP32 und Pi.

\item[CRC]
Cyclic Redundancy Check – Prüfsumme für Datenintegrität.\\
\textit{Im Projekt:} CRC-16 CCITT für Serial-Frames.

\end{description}

% -----------------------------------------------------------------------------
\subsection*{D}

\begin{description}[style=nextline, leftmargin=2.5cm, labelwidth=2.3cm]

\item[Debounce]
Entprellen von mechanischen Tastern.\\
\textit{Im Projekt:} \SI{50}{\milli\second} Debounce-Zeit, Edge-basiert.

\item[Datenfluss]
Verarbeitungssequenz: \code{read → update → write}\\
\textit{Im Projekt:} CD4021 → Logic → HC595

\end{description}

% -----------------------------------------------------------------------------
\subsection*{E}

\begin{description}[style=nextline, leftmargin=2.5cm, labelwidth=2.3cm]

\item[Edge]
Flanke: Übergang von einem Zustand zum anderen.
\begin{itemize}[nosep, leftmargin=1.5em]
\item \textbf{Rising Edge:} 0 → 1
\item \textbf{Falling Edge:} 1 → 0
\end{itemize}
\textit{Im Projekt:} Press-Edge = Taste wird gedrückt.

\item[E2E-Test]
End-to-End Test mit echter Hardware.\\
\textit{Im Projekt:} \code{test\_e2e\_hw} prüft SPI, Timing, Mapping.

\end{description}

% -----------------------------------------------------------------------------
\subsection*{F}

\begin{description}[style=nextline, leftmargin=2.5cm, labelwidth=2.3cm]

\item[FreeRTOS]
Real-Time Operating System für Mikrocontroller.\\
\textit{Im Projekt:} Optional, für Task-basierte Architektur.

\item[Frame]
Datenpaket mit Header, Payload, CRC.\\
\textit{Im Projekt:} COBS-kodiert über Serial.

\end{description}

% -----------------------------------------------------------------------------
\subsection*{G}

\begin{description}[style=nextline, leftmargin=2.5cm, labelwidth=2.3cm]

\item[GPIO]
General Purpose Input/Output – universelle Pins.

\item[Guard]
Zeitfenster, das erneute Aktionen blockiert.\\
\textit{Im Projekt:} Press-Guard verhindert Doppel-Klicks.

\end{description}

% -----------------------------------------------------------------------------
\subsection*{H}

\begin{description}[style=nextline, leftmargin=2.5cm, labelwidth=2.3cm]

\item[HAL]
Hardware Abstraction Layer – Hardware-Abstraktion.\\
\textit{Im Projekt:} \code{src/hal/} für SPI-Zugriff.

\item[HC595]
74HC595: 8-Bit Schieberegister mit Latch.\\
\textit{Im Projekt:} Steuert LEDs (Active-High).

\end{description}

% -----------------------------------------------------------------------------
\subsection*{I}

\begin{description}[style=nextline, leftmargin=2.5cm, labelwidth=2.3cm]

\item[ISR]
Interrupt Service Routine – Interrupt-Handler.\\
\textit{Regel:} Kurz halten, keine langen Operationen.

\end{description}

% -----------------------------------------------------------------------------
\subsection*{L}

\begin{description}[style=nextline, leftmargin=2.5cm, labelwidth=2.3cm]

\item[Latch]
Speichert Daten für stabile Ausgabe.\\
\textit{Im Projekt:} HC595 RCK-Pin triggert Latch.

\item[LSB / MSB]
\begin{itemize}[nosep, leftmargin=1.5em]
\item \textbf{LSB:} Least Significant Bit (niedrigstes Bit)
\item \textbf{MSB:} Most Significant Bit (höchstes Bit)
\end{itemize}
\textit{Im Projekt:} MSB-first für Buttons und LEDs.

\end{description}

% -----------------------------------------------------------------------------
\subsection*{M}

\begin{description}[style=nextline, leftmargin=2.5cm, labelwidth=2.3cm]

\item[Maske]
Bitmuster zum Selektieren/Filtern von Bits.
\begin{lstlisting}[style=arduino,numbers=none]
value & mask;   // Bits auslesen
value |= mask;  // Bits setzen
value &= ~mask; // Bits loeschen
\end{lstlisting}

\item[MISO / MOSI]
\begin{itemize}[nosep, leftmargin=1.5em]
\item \textbf{MISO:} Master In, Slave Out (Daten zum Master)
\item \textbf{MOSI:} Master Out, Slave In (Daten zum Slave)
\end{itemize}

\end{description}

% -----------------------------------------------------------------------------
\subsection*{O}

\begin{description}[style=nextline, leftmargin=2.5cm, labelwidth=2.3cm]

\item[OE]
Output Enable – Aktiviert Ausgänge (meist active-low).\\
\textit{Im Projekt:} HC595 OE = LOW → LEDs aktiv.

\item[One-Hot]
Genau ein Bit gesetzt (oder keins).\\
\textit{Im Projekt:} Immer nur eine LED aktiv.

\end{description}

% -----------------------------------------------------------------------------
\subsection*{P}

\begin{description}[style=nextline, leftmargin=2.5cm, labelwidth=2.3cm]

\item[Pipeline]
Datenfluss durch mehrere Verarbeitungsstufen.\\
\textit{Im Projekt:} read → debounce → selection → write

\item[Polling]
Regelmäßiges Abfragen eines Zustands.\\
\textit{Im Projekt:} \SI{10}{\milli\second} Scan-Intervall.

\item[Pull-up / Pull-down]
Widerstand, der Pin auf definierten Pegel zieht.\\
\textit{Im Projekt:} Pull-ups für CD4021-Inputs.

\end{description}

% -----------------------------------------------------------------------------
\subsection*{R}

\begin{description}[style=nextline, leftmargin=2.5cm, labelwidth=2.3cm]

\item[RTOS]
Real-Time Operating System.\\
\textit{Im Projekt:} Optional FreeRTOS.

\end{description}

% -----------------------------------------------------------------------------
\subsection*{S}

\begin{description}[style=nextline, leftmargin=2.5cm, labelwidth=2.3cm]

\item[Scan-Zyklus]
Ein Durchlauf: Lesen → Verarbeiten → Schreiben.

\item[SPI]
Serial Peripheral Interface – synchroner Bus.\\
\textit{Im Projekt:}
\begin{itemize}[nosep, leftmargin=1.5em]
\item CD4021: Mode 1, \SI{500}{\kilo\hertz}
\item HC595: Mode 0, \SI{1}{\mega\hertz}
\end{itemize}

\item[SSOT]
Single Source of Truth – eine Quelle für Daten.\\
\textit{Im Projekt:} Hardware-Parameter nur in \code{include/hw/}.

\end{description}

% -----------------------------------------------------------------------------
\subsection*{T}

\begin{description}[style=nextline, leftmargin=2.5cm, labelwidth=2.3cm]

\item[Task]
Ausführungseinheit in FreeRTOS.\\
\textit{Im Projekt:} \code{io\_task} für IO-Verarbeitung.

\item[Timing]
Zeitliche Abfolge von Signalen.\\
\textit{Kritisch:} Zu schnell/langsam → Fehler.

\end{description}

% -----------------------------------------------------------------------------
\subsection*{U}

\begin{description}[style=nextline, leftmargin=2.5cm, labelwidth=2.3cm]

\item[Unit-Test]
Test einer einzelnen Komponente.\\
\textit{Im Projekt:} Tests für Debounce, Selection, Pipeline.

\item[USB-CDC]
USB Communication Device Class – virtuelle serielle Schnittstelle.

\end{description}

% -----------------------------------------------------------------------------
\subsection*{V}

\begin{description}[style=nextline, leftmargin=2.5cm, labelwidth=2.3cm]

\item[Volatile]
C-Keyword: Variable kann sich unerwartet ändern (ISR, Hardware).
\begin{lstlisting}[style=arduino,numbers=none]
volatile uint8_t flag;
\end{lstlisting}

\end{description}

% -----------------------------------------------------------------------------
\subsection*{W}

\begin{description}[style=nextline, leftmargin=2.5cm, labelwidth=2.3cm]

\item[Watchdog]
Timer, der bei Timeout einen Reset auslöst.\\
\textit{Zweck:} Erkennung von Hängern.

\item[Wire Frame]
Datenpaket auf dem Draht (nach Encoding).

\end{description}

% -----------------------------------------------------------------------------
\subsection*{Weiterführende Dokumentation}

\begin{description}
\item[\code{firmware/docs/overview.md}] Projektübersicht
\item[\code{firmware/docs/wiring.md}] Hardware-Verdrahtung
\item[\code{docs/Serial-Protokoll.md}] Protokoll-Details
\end{description}
